% Section 4.3 - Analisis Hasil Ablation Study
% Untuk dimasukkan ke dalam paper utama

\subsection{Studi Ablasi: Pengaruh Text-Guided Discriminator}

Studi ablasi dilakukan untuk memvalidasi kontribusi arsitektur dual-modal discriminator terhadap performa restorasi dokumen yang berorientasi pada HTR. Tiga konfigurasi dibandingkan: (1) Single-Modal (image-only discriminator), (2) Dual-Modal dengan ground truth text, dan (3) Dual-Modal dengan predicted text dari recognizer. Semua konfigurasi menggunakan generator yang sama (Enhanced) dan dilatih selama 10 epoch dengan dataset sintetis yang sama.

\subsubsection{Temuan Kritis: PSNR Tinggi Tidak Menjamin Readability}

Hasil ablasi mengungkap temuan penting yang mengubah pemahaman mengenai metrik evaluasi untuk document restoration task (Tabel~\ref{tab:ablation_results}). Konfigurasi single-modal mencapai PSNR sebesar 20.01 dB dan SSIM 0.9217, yang mengindikasikan kualitas visual yang baik secara pixel-level. Namun, ketika dievaluasi menggunakan HTR recognizer, konfigurasi ini \textbf{gagal total} dengan Character Error Rate (CER) mencapai 1.0000 (100\% kesalahan). Artinya, tidak ada satu karakter pun yang dapat dikenali dengan benar oleh model HTR.

Fenomena ini terjadi karena discriminator image-only hanya mengevaluasi kualitas visual pixel-level tanpa mempertimbangkan struktur karakter yang penting untuk readability. Generator yang dioptimasi dengan adversarial loss dari discriminator semacam ini cenderung menghasilkan citra yang ``terlihat bersih'' secara visual, namun dengan struktur stroke yang terputus, jarak antar karakter yang tidak konsisten, atau bentuk karakter yang terdistorsi---semua karakteristik yang tidak tertangkap oleh metrik PSNR namun fatal untuk HTR.

\subsubsection{Superioritas Dual-Modal: Peningkatan Readability 59.1\%}

Konfigurasi dual-modal dengan ground truth text menunjukkan performa superior pada semua metrik. Dibandingkan dengan single-modal, konfigurasi ini mencapai:

\begin{itemize}
    \item \textbf{PSNR:} 20.23 dB (+0.22 dB atau +1.1\%)---peningkatan kualitas visual
    \item \textbf{SSIM:} 0.9245 (+0.0028 atau +0.3\%)---peningkatan structural similarity
    \item \textbf{CER:} 0.4092 (-0.5908 atau \textbf{-59.1\%})---peningkatan readability yang sangat signifikan
    \item \textbf{WER:} 0.9322 (-0.0678 atau -6.8\%)---peningkatan word-level recognition
\end{itemize}

Peningkatan CER sebesar 59.1\% (dari 1.0000 menjadi 0.4092) adalah temuan paling krusial dari studi ini. Angka tersebut menunjukkan bahwa \textit{text-guided adversarial training} tidak hanya memberikan perbaikan inkremental, namun mengubah output dari ``tidak dapat dibaca sama sekali'' menjadi ``dapat dibaca dengan tingkat kesalahan 40\%''. Validasi statistik menggunakan uji hipotesis satu sisi mengkonfirmasi signifikansi perbaikan ini ($p < 0.001$).

\subsubsection{Mekanisme Text-Guided Training}

Keunggulan dual-modal discriminator dapat dijelaskan melalui tiga mekanisme utama:

\begin{enumerate}
    \item \textbf{CTC Loss (bobot 10.0):} Memberikan supervisi eksplisit pada tingkat karakter, memaksa generator untuk mempertahankan struktur karakter yang dapat dikenali oleh HTR recognizer. Loss ini di-clip pada nilai 300 untuk stabilitas training.
    
    \item \textbf{Text LSTM dalam Discriminator:} Mengevaluasi sekuens text features yang diekstrak dari citra, memberikan feedback mengenai koherensi tekstual dan readability. Pathway ini memiliki dimensi 512 yang dibalance dengan image features (1024-dim) pada rasio 2:1.
    
    \item \textbf{RecFeat Loss:} Menyelaraskan fitur intermediate antara citra terdegradasi dan citra hasil restorasi pada representasi recognizer, mendorong generator untuk menghasilkan output yang ``recognizable'' bagi model HTR.
\end{enumerate}

Kombinasi ketiga mekanisme ini menciptakan signal pembelajaran yang multi-objektif: generator tidak hanya dioptimasi untuk kualitas visual (pixel-level fidelity), namun juga untuk readability (character-level structure preservation). Hasil ablasi membuktikan bahwa kedua objektif ini dapat dicapai secara bersamaan tanpa trade-off signifikan---bahkan terjadi peningkatan pada kedua aspek.

\subsubsection{Konfigurasi Dual-Modal dengan Predicted Text}

Konfigurasi ketiga menggunakan predicted text dari recognizer (bukan ground truth) untuk mensimulasikan skenario real-world deployment dimana ground truth text tidak selalu tersedia. Hasil menunjukkan performa yang masih kompetitif:

\begin{itemize}
    \item PSNR: 20.07 dB (hampir setara dengan single-modal)
    \item CER: 0.4276 (peningkatan 57.2\% dibanding single-modal)
\end{itemize}

Meskipun CER sedikit lebih tinggi dibanding konfigurasi GT (0.4276 vs 0.4092), peningkatan relatif terhadap single-modal tetap massive (57.2\%). Hal ini mengindikasikan bahwa text-guided training tetap efektif meskipun menggunakan predicted text yang mungkin mengandung error, menjadikan pendekatan ini praktis untuk aplikasi real-world.

\subsubsection{Implikasi untuk Research dan Praktik}

Temuan studi ablasi ini memiliki beberapa implikasi penting:

\begin{enumerate}
    \item \textbf{Metrik Evaluasi:} Untuk document restoration yang berorientasi pada HTR downstream task, metrik visual seperti PSNR dan SSIM tidak cukup. CER/WER harus dilaporkan sebagai metrik primer untuk menilai keberhasilan restorasi.
    
    \item \textbf{Arsitektur Design:} Text-guided discriminator dengan balanced image-text features terbukti esensial, bukan optional. Single-modal approach menghasilkan output yang tidak dapat digunakan untuk HTR aplikasi.
    
    \item \textbf{Multi-Objective Optimization:} GAN framework memungkinkan optimasi multi-objektif (visual quality + readability) secara bersamaan melalui carefully designed discriminator architecture dan loss function combination.
    
    \item \textbf{Generalisasi:} Prinsip dual-modal approach dapat digeneralisasi untuk task lain yang memerlukan visual quality dan task-specific performance (contoh: OCR untuk printed documents, scene text recognition, dll).
\end{enumerate}

Studi ablasi ini memberikan bukti empiris kuat bahwa novelty claim penelitian ini---peningkatan readability melalui dual-modal discriminator---terbukti valid secara kuantitatif dan statistik. Peningkatan CER sebesar 59.1\% dengan minimal trade-off PSNR (+0.22 dB) menunjukkan bahwa pendekatan ini memberikan kontribusi signifikan untuk state-of-the-art dalam document restoration research.
